\documentclass[11pt,a4paper]{article}
\usepackage[margin=1in]{geometry}
\usepackage{hyperref}
\usepackage{enumitem}
\usepackage{graphicx}
\usepackage{xcolor}

% -----------------------------------------------------
% bvraghav-tiet-ucs503p-article.sty
% -----------------------------------------------------

% Variable: Lab Instructor
\makeatletter \newcommand{\BvrSetLabInstructor}[1]{%
  \gdef\Bvr@LabInstructor{#1}}
% default
\newcommand{\Bvr@LabInstructor}{Dr. Raghav %
  B. Venkataramaiyer}
% public accessor
\newcommand{\BvrLabInstructorValue}{\Bvr@LabInstructor}
\makeatother

% Add subtitle
\usepackage{titling}
% -- Set title bold
\pretitle{\begin{center}\bfseries\fontsize{18bp}{18bp}\selectfont}
\makeatletter
\newcommand{\BvrSubtitle}[1]{%
  \posttitle{%
    \par\end{center}
    \begin{center}\large#1\end{center}
    \vskip 0.5em}%
}
\makeatother

% Hints in the section title(s)
\newcommand{\BvrHint}[1]{%
  {\normalfont\normalsize\itshape\hfill #1}}

% -----------------------------------------------------

\title{Grievance-Resolution System for a Tier-2/3 City}

\BvrSetLabInstructor{Dr. Raghav B. Venkataramaiyer}

\BvrSubtitle{%
  Project Proposal for UCS503P  \\
  Submitted to: \BvrLabInstructorValue \\ \bfseries
  Thapar Institute of Engineering and Technology% 
}

\author{
  Vidya Vaidyanathan, CSED%
  \thanks{Roll No: \texttt{1024030xxx}, %
    Email: \texttt{vvaidhanathan\_be24\@thapar.edu}}
  \and 
  Bhanurekha Ganesan, CSED%
  \thanks{Roll No: \texttt{1024030xxx}, %
    Email: \texttt{bganesan\_be24\@thapar.edu}}
}

\date{\today}

\begin{document}
\maketitle

\section{Higher-order goal%
  \BvrHint{\ldots secondary framing}}
This project aims to improve sustainable urban
governance by reducing the time and effort required to
resolve citizen issues. While the topic is broader than
any single municipality, the immediate engineering
objective is to translate \emph{efficient public
  service delivery} into a measurable, deployable
software solution.

\section{Time-to-value %
\BvrHint{\ldots primary heuristic}}
\begin{itemize}[leftmargin=0.7cm]
    \item We will deliver meaningful increments quickly by prototyping the core user workflow and validating it with users within a short iteration window.
    \item We will prioritize fast feedback over perfection by using weekly iterations (and targeting CI-backed automation early) to reduce release friction.
    \item Success is defined by what can be delivered and measured in the first iteration, then improved based on user feedback.
\end{itemize}

\section{Problem Statement}
Citizens in a tier-2/3 Indian city face delays and uncertainty when reporting local issues (e.g., streetlight outages, water leakage, garbage accumulation, road potholes). Common pain points include:
\begin{itemize}[leftmargin=0.7cm]
    \item \textbf{Fragmentation:} complaints are often made via multiple channels (calls, social media, paper forms), causing lost context.
    \item \textbf{Low transparency:} citizens cannot easily see the status of their complaint (e.g., ``received'', ``assigned'', ``in progress'', ``resolved'').
    \item \textbf{Slow coordination:} departments/field teams lack a unified, trackable workflow with clear ownership.
    \item \textbf{Poor data quality:} duplicate or incomplete reports force rework and extend resolution time.
\end{itemize}

\textbf{Why this matters now (contextual relevance):} In an urban setting with limited staff bandwidth, even small process improvements can noticeably reduce resolution delays and repeated submissions.

\section{Proposed Solution %
\BvrHint{\ldots Software Proposition}}
\subsection{Overview}
Build a \textbf{web-based} ``Grievance-to-Resolution'' system with the following components:
\begin{itemize}[leftmargin=0.7cm]
    \item \textbf{Citizen submission portal:} users submit an issue with photo (optional), category, and locality/ward.
    \item \textbf{Automated triage:} basic duplicate detection and completeness checks.
    \item \textbf{Department workflow:} assign issues to relevant department/team and update statuses.
    \item \textbf{Status transparency:} citizens track status and receive notifications when resolved.
    \item \textbf{Admin/ops dashboard:} performance summaries (resolution time, backlog).
\end{itemize}

\subsection{Core workflow}
\begin{enumerate}[leftmargin=0.7cm]
    \item Citizen creates a complaint and receives a complaint ID immediately.
    \item System validates completeness and suggests possible duplicates (lightweight similarity on text/category and approximate location).
    \item Department staff update status (Received $\rightarrow$ Assigned $\rightarrow$ In Progress $\rightarrow$ Resolved).
    \item Citizen views status updates and resolution notes.
\end{enumerate}

\subsection{Operational constraints for an educational, tier-2/3 setting}
\begin{itemize}[leftmargin=0.7cm]
    \item Keep the solution usable on low-to-mid range devices via responsive web design.
    \item Avoid dependence on advanced research algorithms; focus on workflow, automation, and system correctness.
    \item Design for deployability using common web hosting and managed databases.
\end{itemize}

\section{Solution Approach %
\BvrHint{\ldots Engineering focus}}
This is an engineering course project, so we focus on scalable administrative and workflow aspects rather than novel algorithm research.

\subsection{Duplicate/quality assistance %
\BvrHint{\ldots non-research}}
Duplicate detection will be heuristic-based:
\begin{itemize}[leftmargin=0.7cm]
    \item Normalize text (category + short description keywords).
    \item Use a simple similarity measure (e.g., token overlap / TF-IDF cosine with threshold) without claiming state-of-the-art.
    \item Show top candidate duplicates to staff for decision; never fully auto-delete or auto-reject.
\end{itemize}

\subsection{Web architecture %
\BvrHint{\ldots web-based solution}}
A typical 3-tier web architecture:
\begin{itemize}[leftmargin=0.7cm]
    \item \textbf{Frontend:} responsive UI for citizens and staff (status pages, submission form, dashboard).
    \item \textbf{Backend API:} authentication, complaint submission, status transitions, assignment logic.
    \item \textbf{Database:} store users, complaints, status history, attachments metadata.
\end{itemize}

\subsection{CI/CD and fast delivery enabler}
CI/CD will be used to:
\begin{itemize}[leftmargin=0.7cm]
    \item Automatically build and run tests on every change.
    \item Produce repeatable deployment artifacts to staging.
    \item Enable safe and frequent releases of incremental features.
\end{itemize}

\section{Evaluation Criterion %
\BvrHint{\ldots measurable and attributable}}
We define success using measurable metrics tied to the core workflow.

\subsection{Primary evaluation metric}
\textbf{Time-to-acknowledgement (TTA):} time between complaint submission and ``Received/Assigned'' status creation in the system.
\begin{itemize}[leftmargin=0.7cm]
    \item Target: median TTA $\le$ 2 hours during pilot window.
    \item Attribution: measured from database timestamps (submission vs. first status update by staff).
\end{itemize}

\subsection{Secondary evaluation metrics}
\begin{itemize}[leftmargin=0.7cm]
    \item \textbf{Resolution time reduction:} reduce time-to-resolved compared to baseline process (even if baseline is collected via survey).
    \item \textbf{Duplicate rate:} percentage of complaints flagged as duplicates and resolved without rework.
    \item \textbf{User clarity:} citizen-reported clarity of status using a simple 1--5 feedback question.
    \item \textbf{Reliability:} availability of staging/production for login and submission during business hours (e.g., $\ge$ 99\% uptime in pilot).
\end{itemize}

\subsection{Pilot validation plan %
\BvrHint{\ldots high-level}}
\begin{itemize}[leftmargin=0.7cm]
    \item Recruit 10--20 pilot users (citizens) and 3--5 staff roles (department simulation if real staff unavailable).
    \item Compare metrics before/after within the iteration window by logging baseline estimates via short interviews or existing process observations.
\end{itemize}

\section{Scalability %
\BvrHint{\ldots with foresight}}
The proposition should scale in both theoretical and practical senses.

\begin{enumerate}[leftmargin=0.7cm]
    \item \textbf{Scalable (theoretically):} The system should support growth in number of complaints and concurrent dashboard usage by:
    \begin{itemize}[leftmargin=0.7cm]
        \item decoupling components (frontend/backend),
        \item enabling pagination and indexing for common queries,
        \item using stateless API design.
    \end{itemize}
    
    \item \textbf{Operationally deployable (practically):} The system should run on common web hosting:
    \begin{itemize}[leftmargin=0.7cm]
        \item managed database,
        \item containerized or platform-hosted deployment,
        \item CI/CD pipeline for staging/production.
    \end{itemize}
\end{enumerate}

\section{Engine availability heuristic}
A mature set of ``engines'' (tooling/services) is available to implement this as a web-based project:
\begin{itemize}[leftmargin=0.7cm]
    \item Web framework for REST APIs and templated/reactive pages.
    \item Managed relational or document database.
    \item Authentication approach (token-based or session-based).
    \item Object storage for optional attachments (or local storage for pilot).
    \item Test framework for unit/integration tests.
    \item CI runner and staging deployment target.
\end{itemize}
We will use these mature components to focus on workflow design, correctness, and measurement rather than inventing new infrastructure.

\section{Project scope and deliverables %
\BvrHint{\ldots iteration-friendly}}
\subsection{Initial deliverable %
\BvrHint{\ldots first iteration}}
\begin{itemize}[leftmargin=0.7cm]
    \item Citizen submission form + complaint ID generation
    \item Staff workflow: update status and assignment
    \item Citizen status page: view timeline/history
    \item Basic logging and timestamps for TTA measurement
    \item CI: build + backend unit tests + linting
    \item CD to staging with a single command or pipeline job
\end{itemize}

\subsection{Subsequent deliverables}
\begin{itemize}[leftmargin=0.7cm]
    \item Lightweight duplicate suggestions with staff review
    \item Notification layer (email/SMS optional in later iteration; can start with in-app notifications)
    \item Admin dashboard for metrics and backlog
    \item Improved search/filtering and indexed queries for scale readiness
\end{itemize}

\section{Risks and mitigations}
\begin{itemize}[leftmargin=0.7cm]
    \item \textbf{Data privacy / consent:} minimize collected data; keep attachments optional; store only what is needed.
    \item \textbf{Staff adoption:} provide an easy UI for status updates; keep workflow steps minimal; include shortcuts.
    \item \textbf{Evaluation measurement ambiguity:} instrument timestamps and define clear status transitions before development begins.
    \item \textbf{Operational issues in pilot:} use staging early; ensure CI/CD produces repeatable deployments.
\end{itemize}

\section{Summary}
This project proposes a deployable web platform to reduce grievance handling delays in a tier-2/3 city context. It meets the course criteria by emphasizing time-to-value through rapid iterations, implementing CI/CD to support frequent safe releases, and using measurable evaluation criteria (especially time-to-acknowledgement). It remains focused on engineering workflow and scalability readiness rather than research-driven novelty.

\end{document}
